\phantomsection
\color{unidarkgreen}\section*{\centering{\large{ПЕРЕЧЕНЬ СОКРАЩЕНИЙ}}}
\addcontentsline{toc}{section}{Перечень сокращений}
{\color{white}\fontsize{0.1pt}{0.1pt}\selectfont<begABBRS>}
\color{black}
\begin{longtable}{>{\raggedright\arraybackslash}m{2cm}>{\raggedright\arraybackslash}m{0.5cm}>{\raggedright\arraybackslash}m{20cm}}
\endfirsthead\endhead\endfoot\endlastfoot
АМТЗ & -- & аварийная максимальная токовая защита; \\
АПВ & -- & автоматическое повторное включение; \\
АУ & -- & автоматическое ускорение; \\
АЦП & -- & аналого-цифровой преобразователь; \\
БИ & -- & блок испытательный; \\
БК & -- & блокировка при качаниях; \\
БНН & -- & блокировка при неисправности цепей напряжения; \\
БНТ & -- & блокировка при бросках намагничивающего тока; \\
БСТО & -- & блокировка при сквозных токах через ошиновку; \\
БЧС & -- & блокировка чувствительных ступеней; \\
ВЛ & -- & воздушная линия; \\
ВН & -- & высшее напряжение; \\
ВО & -- & вторичная обмотка; \\
ВЧ & -- & высокочастотный; \\
ВЧПП & -- & высокочастотный приемо-передатчик; \\
ВЧТО & -- & высокочастотное телеотключение; \\
ДЗ & -- & дистанционная защита; \\
ДЗШ & -- & дифференциальная защита шин; \\
ЕЭС & -- & единая энергетическая система; \\
ЗИП & -- & запасные части, инструменты и принадлежности; \\
ЗНР & -- & защита от неполнофазного режима; \\
ЗНФ & -- & защита от непереключения фаз; \\
ИО & -- & измерительный орган / избирательный орган; \\
ИЧМ & -- & интерфейс человек-машина; \\
КЗ & -- & короткое замыкание; \\
КОН & -- & контроль отсутствия напряжения; \\
КС & -- & контроль синхронизма / контрольная сумма; \\
КСВ & -- & контроль силового выключателя; \\
КСН & -- & контроль синхронизма и напряжения; \\
ЛВ & -- & логика включения; \\
ЛО & -- & логика отключения; \\
ЛОСК & -- & логика отключения слабого конца; \\
ЛР & -- & линейный разъединитель; \\
ЛС & -- & логика связи; \\
ЛЭП & -- & линия электропередачи; \\
МФТО & -- & междуфазная токовая отсечка; \\
МЭК & -- & международная электротехническая комиссия; \\
НВЧЗ & -- & направленная высокочастотная защита; \\
НЗ & -- & нормально закрытый; \\
НН & -- & низшее напряжение; \\
ННЛ & -- & наличие напряжения на линии; \\
ННЭ & -- & наличие напряжения на энергообъекте; \\
НО & -- & нормально открытый; \\
НП & -- & нулевая последовательность; \\
ОВ & -- & оперативный вывод / обходной выключатель; \\
ОН & -- & орган направления / отключение нагрузки; \\
ОНЛ & -- & отсутствие напряжения на линии; \\
ОНМ & -- & орган направления мощности; \\
ОНМНП & -- & орган направления мощности нулевой последовательности; \\
ОНМОП & -- & орган направления мощности обратной последовательности; \\
ОНЭ & -- & отсутствие напряжения на энергообъекте; \\
ООВП & -- & орган определения вида повреждения; \\
ОП & -- & обратная последовательность; \\
ОТ & -- & оперативный ток; \\
ОТФ & -- & отключение трех фаз; \\
ОУ & -- & оперативное ускорение; \\
ПАВ & -- & полуавтоматическое включение; \\
ПО & -- & пусковой орган / программное обеспечение; \\
ПУ & -- & пульт управления / поперечное ускорение; \\
РД & -- & реле давления; \\
РЗ & -- & релейная защита; \\
РЗА & -- & релейная защита и автоматика; \\
РКВ & -- & реле команды «включить»; \\
РКО & -- & реле команды «отключить»; \\
РПВ & -- & реле положения «включено»; \\
РПО & -- & реле положения отключено; \\
РС & -- & регистрация событий / реле сопротивления; \\
РТ & -- & реле тока; \\
РУ & -- & распределительное устройство; \\
РФ & -- & реле фиксации; \\
РФК & -- & реле фиксации команд; \\
РФП & -- & реле фиксации включенного положения; \\
РЭ & -- & руководство по эксплуатации; \\
СО & -- & система охлаждения; \\
СШ & -- & система (секция) шин; \\
ТЗНП & -- & токовая защита нулевой последовательности; \\
ТЗО & -- & токовая защита ошиновки; \\
ТН & -- & трансформатор напряжения; \\
ТНЗНП & -- & токовая направленная защита нулевой последовательности; \\
ТО & -- & токовая отсечка; \\
ТС & -- & технологическая сигнализация; \\
ТТ & -- & трансформатор тока; \\
ТУ & -- & телеуправление / телеускорение / технические условия; \\
УВ & -- & управление выключателем / управляющее воздействие; \\
УКЕТ & -- & устройство компенсации емкостных токов; \\
УРОВ & -- & устройство резервирования при отказе выключателя; \\
УС & -- & улавливание синхронизма; \\
ФК & -- & функциональная клавиша; \\
ФПВ & -- & фиксация положения выключателя; \\
ФСЛ & -- & фиксация состояния ЛЭП; \\
ФСУ & -- & функциональная схема устройства; \\
ЦН & -- & цепи напряжения; \\
ЦОС & -- & цифровая обработка сигналов; \\
ЦП & -- & центральный процессор; \\
ЦУ & -- & цепи управления; \\
ЧС & -- & чувствительная ступень; \\
ШОН & -- & шкаф отбора напряжения; \\
ШП & -- & шинки питания; \\
ШСВ & -- & шиносоединительный выключатель; \\
ШЭТ & -- & шкаф электротехнический типовой; \\
ЭМВ & -- & электромагнит включения; \\
ЭМО & -- & электромагнит отключения; \\
ЭМУ & -- & электромагнит управления. \\
{\color{white}\fontsize{0.1pt}{0.1pt}\selectfont<endABBRS>}
\end{longtable}
